% ^^A -*- japanese-latex -*-
%
% \ifx\epTeXinputencoding\undefined\else
%   \epTeXinputencoding utf8
% \fi
% \iffalse
%<*driver>
\ifx\epTeXinputencoding\undefined\else
  \epTeXinputencoding utf8
\fi
%</driver>
%
%
%<*driver>
\documentclass[dvipdfmx,uplatex,a4paper,openany,jbase=12Q,nomag*,textwidth=40zw%draft
               ]{bxjsarticle}
\setpagelayout*{includefoot,%
                vdivide={20mm,,20mm},%
                hdivide={2in,,1in},%
                footskip=10mm}
%\documentclass[lualatex,a4paper,12Q]{jlreq}
\usepackage{doc}
\usepackage{ifptex}
\ifptex
  \xspcode"5C=1 %% \
  \xspcode"22=1 %% "
\fi
%\addtolength{\textwidth}{-1in}
%\addtolength{\evensidemargin}{1in}
%\addtolength{\oddsidemargin}{1in}
%\addtolength{\marginparwidth}{1in}
\setlength\marginparpush{0pt}
\usepackage{etoolbox}
% \usepackage{listings}
% \ifptex
%   \usepackage{plistings}
% \fi
\usepackage{color}
\usepackage{xspace}
\usepackage{ifpdf}
% \usepackage[colorlinks=true]{hyperref}
% \ifptex
%   \usepackage{pxjahyper}
% \fi
\usepackage[T1]{fontenc}
\usepackage{lmodern}
\definecolor{myblue}{rgb}{0,0,0.75}
\definecolor{mygreen}{rgb}{0,0.45,0}
\CodelineNumbered
\EnableCrossrefs
\CodelineIndex
\RecordChanges % Gather update information
\setcounter{StandardModuleDepth}{1}
\setcounter{tocdepth}{3}

\def\TikZ{Ti\textsl{k}Z}
\newcommand*{\Lcount}[1]{\textsl{\small#1}}
\newcommand*{\Lopt}[1]{\textsf{#1}}
\newcounter{@clineno}
\def\mlineplus#1{\setcounter{@clineno}{\arabic{CodelineNo}}%
   \addtocounter{@clineno}{#1}\arabic{@clineno}}
\makeatletter
\DeclareRobustCommand{\cmd}[1]{\cs{\expandafter\cmd@to@cs\string#1}}
\def\cmd@to@cs#1#2{\char\number`#2\relax}
\makeatother
\DeclareRobustCommand{\cs}[1]{\texttt{\char`\\#1}}
\providecommand{\marg}[1]{%
  {\ttfamily\char`\{}\meta{#1}{\ttfamily\char`\}}}
\providecommand{\argchar}[1]{%
  {\ttfamily\char`\{}\texttt{#1}{\ttfamily\char`\}}}
\providecommand{\oarg}[1]{%
  {\ttfamily[}\meta{#1}{\ttfamily]}}
\newrobustcmd*{\env}[1]{\mbox{\ttfamily#1}}
\let\cntname\env
\newrobustcmd*{\file}[1]{\mbox{\ttfamily#1}}
\newrobustcmd*{\prm}[1]{%
  \ifblank{#1}%
    {}%
    {\mbox{%
       \ensuremath\langle
       \normalfont\textit{#1}%
       \ensuremath\rangle}}}

\usepackage{bookkeeping}
\hypersetup{colorlinks,hyperfootnotes=false,linkcolor=myblue,urlcolor=mygreen}

\GetFileInfo{bookkeeping.sty}
\begin{document}
  \DocInput{bookkeeping.dtx}
\end{document}
%</driver>
%
% \fi
% \changes{v1.1.0}{2023/04/26}
%   {make this .dtx file}
% \changes{v1.0.0}{2023/02/05}
%   {make .sty file}
% \DoNotIndex{\@,\@@par,\@beginparpenalty,\@empty}
% \DoNotIndex{\@flushglue,\@gobble,\@input}
% \DoNotIndex{\@makefnmark,\@makeother,\@maketitle}
% \DoNotIndex{\@namedef,\@ne,\@spaces,\@tempa}
% \DoNotIndex{\@tempb,\@tempswafalse,\@tempswatrue}
% \DoNotIndex{\@thanks,\@thefnmark,\@topnum}
% \DoNotIndex{\@@,\@elt,\@forloop,\@fortmp,\@gtempa,\@totalleftmargin}
% \DoNotIndex{\",\/,\@ifundefined,\@nil,\@verbatim,\@vobeyspaces}
% \DoNotIndex{\|,\~,\ ,\active,\advance,\aftergroup,\begingroup,\bgroup}
% \DoNotIndex{\mathcal,\csname,\def,\documentstyle,\dospecials,\edef}
% \DoNotIndex{\egroup}
% \DoNotIndex{\else,\endcsname,\endgroup,\endinput,\endtrivlist}
% \DoNotIndex{\expandafter,\fi,\fnsymbol,\futurelet,\gdef,\global}
% \DoNotIndex{\hbox,\hss,\if,\if@inlabel,\if@tempswa,\if@twocolumn}
% \DoNotIndex{\ifcase,\iftrue,\iffalse}
% \DoNotIndex{\ifcat,\iffalse,\ifx,\ignorespaces,\index,\input,\item}
% \DoNotIndex{\jobname,\kern,\leavevmode,\leftskip,\let,\llap,\lower}
% \DoNotIndex{\m@ne,\next,\newpage,\nobreak,\noexpand,\nonfrenchspacing}
% \DoNotIndex{\obeylines,\or,\protect,\raggedleft,\rightskip,\rm,\sc}
% \DoNotIndex{\setbox,\setcounter,\small,\space,\string,\strut}
% \DoNotIndex{\strutbox}
% \DoNotIndex{\thefootnote,\thispagestyle,\topmargin,\trivlist,\tt}
% \DoNotIndex{\twocolumn,\typeout,\vss,\vtop,\xdef,\z@}
% \DoNotIndex{\,,\@bsphack,\@esphack,\@noligs,\@vobeyspaces,\@xverbatim}
% \DoNotIndex{\`,\catcode,\end,\escapechar,\frenchspacing,\glossary}
% \DoNotIndex{\hangindent,\hfil,\hfill,\hskip,\hspace,\ht,\it,\langle}
% \DoNotIndex{\leaders,\long,\makelabel,\marginpar,\markboth,\mathcode}
% \DoNotIndex{\mathsurround,\mbox,\newcount,\newdimen,\newskip}
% \DoNotIndex{\nopagebreak}
% \DoNotIndex{\parfillskip,\parindent,\parskip,\penalty,\raise,\rangle}
% \DoNotIndex{\section,\setlength,\TeX,\topsep,\underline,\unskip,\verb}
% \DoNotIndex{\vskip,\vspace,\widetilde,\\,\%,\@date,\@defpar}
% \DoNotIndex{\[,\{,\},\]}
% \DoNotIndex{\count@,\ifnum,\loop,\today,\uppercase,\uccode}
% \DoNotIndex{\baselineskip,\begin,\tw@}
% \DoNotIndex{\a,\b,\c,\d,\e,\f,\g,\h,\i,\j,\k,\l,\m,\n,\o,\p,\q}
% \DoNotIndex{\r,\s,\t,\u,\v,\w,\x,\y,\z,\A,\B,\C,\D,\E,\F,\G,\H}
% \DoNotIndex{\I,\J,\K,\L,\M,\N,\O,\P,\Q,\R,\S,\T,\U,\V,\W,\X,\Y,\Z}
% \DoNotIndex{\1,\2,\3,\4,\5,\6,\7,\8,\9,\0}
% \DoNotIndex{\!,\#,\$,\&,\',\(,\),\+,\.,\:,\;,\<,\=,\>,\?,\_}
% \DoNotIndex{\discretionary,\immediate,\makeatletter,\makeatother}
% \DoNotIndex{\meaning,\newenvironment,\par,\relax,\renewenvironment}
% \DoNotIndex{\repeat,\scriptsize,\selectfont,\the,\undefined}
% \DoNotIndex{\arabic,\do,\makeindex,\null,\number,\show,\write,\@ehc}
% \DoNotIndex{\@author,\@ehc,\@ifstar,\@sanitize,\@title,\everypar}
% \DoNotIndex{\if@minipage,\if@restonecol,\ifeof,\ifmmode}
% \DoNotIndex{\lccode,\newtoks,\onecolumn,\openin,\p@,\SelfDocumenting}
% \DoNotIndex{\settowidth,\@resetonecoltrue,\@resetonecolfalse,\bf}
% \DoNotIndex{\clearpage,\closein,\lowercase,\@inlabelfalse}
% \DoNotIndex{\selectfont,\mathcode,\newmathalphabet,\rmdefault}
% \DoNotIndex{\makebox}
% \DoNotIndex{\bfdefault,\centering,\newcommand,\NeedsTeXFormat}
% \DoNotIndex{\hline,\ref,\label,\newcounter,\refstepcounter,\stepcounter}
% \DoNotIndex{\subsection}
% \DoNotIndex{\RequirePackage,\ProvidePackage}
% \DoNotIndex{\numexpr,\dimexpr}
% \DoNotIndex{\zw}
% \DoNotIndex{\multicolumn}
% \iffalse etoolbox \fi
% \DoNotIndex{\csuse,\csdef,\csedef,\csgdef,\csxdef,\letcs}
% \iffalse hyperref \fi
% \DoNotIndex{\hyperref}
% \title{\filename}
% \author{bell}
% \date{\filedate}
% \maketitle
% \tableofcontents
%
% \MakeShortVerb{\|}
%
% \StopEventually{}
% \section{概略・使い方}
% \subsection{勘定科目の登録}
% \cmd{\registerAccountID}\marg{AccType}\marg{AccLegbal}\marg{勘定科目名}で
% 勘定科目を登録します．
% \begin{itemize}
% \item \marg{AccType}は勘定科目の種類で
%   \begin{itemize}
%     \item A: 資産 (Asset)
%     \item L: 負債 (Liability)
%     \item N: 総資産 (Net asset)
%     \item R: 収益 (Revenue)
%     \item E: 費用 (Expence)
%     \item M: その他 (Misc.)
%   \end{itemize}
%   です
% \item \marg{AccLegbal}は借方残高(D)か貸方残高(C)かです．
% \end{itemize}
% \subsection{仕訳の登録}
% \cmd{\journalentry}\marg{月}\marg{日}\marg{借方勘定科目}\marg{金額}\marg{貸方勘定科目}\marg{摘要}で
% 仕訳を登録します．
% \cmd{\journalentry*}\marg{月}\marg{日}\marg{借方勘定科目}\marg{金額}\marg{貸方勘定科目}で
% 摘要なしの仕訳を登録出来ます．
% \begin{itemize}
% \item 今のところ単一仕訳方式のみ対応です．
% \item 開始仕訳はそれ用に勘定科目を設定して行ってください．
% \end{itemize}
% 月替わりには\cmd{\ledgermonthlyclosing}\marg{closing month}を入れます．
% 
% \subsection{調簿の出力}
% 全部終わったら
%   \cmd{\printjournal}で仕訳帳を，
%   \cmd{\printgeneralledger}で総勘定元帳を，
%   \cmd{\printmonthlytrialbalance}で月ごとの残高試算表を出力します，
% \iffalse
% \fi
% \subsection{決算整理仕訳}
% 決算整理仕訳についての一般論．
% \begin{itemize}
%   \item 現金・預金の残高を実際に確認し，帳簿との過不足がないか確認する
%   \item 固定資産の減価償却を行い，減価償却費として費用計上する
%   \item 決算時点の時価にもとづいて有価証券の評価を調整する
%   \item 追加仕訳を仕訳帳に反映させ，総勘定元帳に転記する
% \end{itemize}
% \subsection{決算仕訳}
% \begin{itemize}
%   \item 費用と収益を損益勘定へ振替
%   \item 損益勘定から当期純利益または純損失を算出
%   \item 当期純利益または純損失を資本金勘定へ振替
%   \item 資産，負債および総資産を閉鎖残高勘定へ振替
% \end{itemize}
% 
% \section{実装}
% \subsection{初期設定}
%    \begin{macrocode}
%<*style>
\NeedsTeXFormat{pLaTeX2e}
\ProvidesPackage{bookkeeping}%
  [2023/02/05 v1.1.0 bookkeeping macros]
\RequirePackage[colorlinks=true]{hyperref}
\RequirePackage{pxjahyper}
\RequirePackage{etoolbox}
\RequirePackage{dcolumn}
%    \end{macrocode}
%    \begin{macrocode}

%%%%%%%%%%%%%%%%%%%%%%%%%%%%%%%%%%%%%%%%%%%%%%%%%%%%%%%%%%%%%%%%%%%%%%%%%%%%%%%%
\newcounter{JECount}
\newcounter{JEIDiterator}
\newcounter{AccountTypeCount}
\newcounter{AccountTypeiterator}
\newcounter{AccountIDiterator}
\newcounter{AccLdgIDiterator}
\newcounter{LedgerEntryBalance}
\newcounter{ClosingMonthCount}
\newcounter{MonthIterator}

\newcounter{tableLineCount}
\def\maxtablelinenumber{30}
\def\BK@closingmonth{13}
\newcount\sum@debit
\newcount\sum@credit
\newcount\BKtemp@countA
\newcount\BKtemp@countB
\newcount\BKtemp@countC
\newcount\BKtemp@countD
%%%%%%%%%%%%%%%%%%%%%%%%%%%%%%%%%%%%%%%%%%%%%%%%%%%%%%%%%%%%%%%%%%%%%%%%%%%%%%%%
%    \end{macrocode}
% \subsection{勘定科目の種類および勘定科目}
% \subsubsection{勘定科目の種類}
% \changes{v1.1.0}{2024/02/21}
%   {勘定科目の種類についてリファクタ}
% 勘定科目の種類が\meta{AccType}である勘定科目の種類について、
% 勘定科目の各要素には\cntname{\cmd{\AccType}\meta{AccType}@\meta{prop.}}で
% アクセスできる．
% \meta{prop.}の種類は次の通り．
% \begin{description}
% \item[\cntname{Name}]
%   勘定科目の種類の名前．
% \item[\cntname{ID}]
%   勘定科目の種類の識別番号(ID)．
% \item[\cntname{AccountCount}]
%   勘定科目の種類に登録された勘定科目数のカウンタ．
% \end{description}
% \iffalse
% \begin{itemize}
% \item
%   \cmd{\csuse}\argchar{BK@AccType\meta{AccTypeID}}で\meta{AccType}を得る．
% \item
%   \cntname{AccountCount\meta{AccType}}カウンタで
%   勘定科目の種類が\meta{AccType}の勘定科目の総登録数を得る．
% \end{itemize}
% \fi
%
% 勘定科目の種類の総数はカウンタ\cntname{AccountTypeCount}に保持する．
%
% 勘定科目の種類の識別番号\meta{AccTypeID}から\meta{AccType}を得るために
% \cmd{\AccType@getAccType}\marg{AccTypeID}を用意する．
%
% \begin{macro}{\BK@registerAccType}\marg{AccType}\marg{AccTypeName}\par
% \changes{v1.1.0}{2023/04/26}
%   {the name of registerAccType is changed to BK@registerAccType.}
% \changes{v1.1.0}{2024/02/21}
%   {勘定科目の種類についてリファクタ．}
% 勘定科目の種類\meta{AccType}を登録する．
% カウンタを確保するので無意味に何度も呼び出さないこと．
%    \begin{macrocode}
%% Account
%% registerAccType: 勘定科目の種類の登録
%% #1: AccType
%% #2: AccTypeName
\def\BK@registerAccType#1#2{%
  \stepcounter{AccountTypeCount}%
  \csgdef{AccType#1@Name}{#2}%
  \csxdef{revAccType\the\c@AccountTypeCount @AccType}{#1}%
  \csxdef{AccType#1@ID}{\theAccountTypeCount}%
%%\csdef{BK@AccType\the\c@AccountTypeCount}{#1}%
  \newcounter{AccType#1@AccountCount}%
}
%    \end{macrocode}
% \end{macro}
%    \begin{macrocode}
\BK@registerAccType{A}{資産}% 資産 (Asset)
\BK@registerAccType{L}{負債}% 負債 (Liability)
\BK@registerAccType{N}{総資産}% 総資産 (Net asset)
\BK@registerAccType{R}{収益}% 収益 (Revenue)
\BK@registerAccType{E}{費用}% 費用 (Expence)
\BK@registerAccType{M}{その他}% その他 (Misc.)

%    \end{macrocode}
% \begin{macro}{\AccType@getAccType}\marg{勘定科目の種類ID}\par
% \changes{v1.1.0}{2024/02/21}
%   {勘定科目の種類についてリファクタしたのに伴ってアクセス用のマクロを設置．}
%  勘定科目の種類IDから勘定科目の種類\meta{AccType}を得る．
%    \begin{macrocode}
\def\AccType@getAccType#1{\csuse{revAccType#1@AccType}}
%    \end{macrocode}
% \end{macro}
%
% \subsubsection{勘定科目}
% 勘定科目のIDを\meta{AccType}\meta{AccID}と表すこととすると，
% 勘定科目の各要素には\cntname{\cmd{\Account}\meta{AccType}\meta{AccID}@\meta{prop.}}で
% アクセスできる．
% \meta{prop.}の種類は次の通り．
% \begin{description}
% \item[\cntname{Name}]
%   勘定科目名を得る．
% \item[\cntname{LedgerCount}]
%   勘定科目に登録された取引登録数．
%   この数字でイテレータの分岐をさせる．
% \item[\cntname{LedgerBalanceType}]
%   勘定科目が借方残高(D)か貸方残高(C)か．
% \item[\cntname{Month}\meta{month}@\meta{month prop.}]
%   \meta{month}で指定した月ごとの各要素にアクセスする．
%   \meta{month prop.}の種類は次の通り．
%   \begin{description}
%   \item[\cntname{Balance}]
%     月ごとの残高．
%   \item[\cntname{Debit}]
%     月ごとの借方合計．
%   \item[\cntname{Credit}]
%     月ごとの貸方合計．
%   \end{description}
% \end{description}
% 
% 勘定科目名\meta{AccName}から\meta{AccType}および\meta{AccID}を得るために
% 次のマクロが用意されている．
% \begin{description}
% \item[\cmd{\Account@getAccID}\marg{AccName}]
%   勘定科目名から\meta{AccID}を得る．
% \item[\cmd{\Account@getAccType}\marg{AccName}]
%   勘定科目名から\meta{AccType}を得る．
% \item[\cmd{\Account@getAccTypeAccID}\marg{AccName}]
%   勘定科目名から\meta{AccType}\meta{AccID}を得る．
% \end{description}
\iffalse
0月は開始仕訳用，13月は閉鎖用．
\fi
% \begin{macro}{\registerAccountID}\marg{Account Type}%
%                                  \marg{Account Ledger Balance}%
%                                  \marg{Account Name}\par
% \changes{v1.1.0}{2024/02/27}
%   {勘定科目の種類についてリファクタしたのに伴ってアクセス用のCSを変更．}
% 勘定科目を登録する．
%    \begin{macrocode}
%% registerAccountID: 勘定科目の登録
%% #1: Account type
%% #2: Account Ledger Balance (Debit(D) or Credit(C))
%% #3: Account Name
\newcommand{\registerAccountID}[3]{%
  \global\advance\csuse{c@AccType#1@AccountCount}\@ne\relax%
  \csgdef{revAccount#3@Type}{#1}% AccType <- Account Name
  \csxdef{revAccount#3@ID}%
    {\the\csname c@AccType#1@AccountCount\endcsname}% ID <- Account Name
  \csgdef{Account#1\the\csuse{c@AccType#1@AccountCount}@Name}{#3}%
  \csgdef{Account#1\the\csuse{c@AccType#1@AccountCount}@LedgerCount}{0}%
  \csgdef{Account#1\the\csuse{c@AccType#1@AccountCount}@LedgerBalanceType}{#2}%
  \BKtemp@countA=-1\relax
  \edef\reserved@a{#1\the\csuse{c@AccType#1@AccountCount}}%
  \register@Acc@Mon@Bal
% \typeout{<AccType>: #1\par <AccID>: \the\csuse{c@AccType#1@AccountCount}}%
}
%    \end{macrocode}
% \end{macro}
% \begin{macro}{\register@Acc@Mon@Bal}\par
% 勘定科目の月ごとの残高・借方合計・貸方合計の登録．
%    \begin{macrocode}
% \register@Acc@Mon@Bal: 勘定科目の月ごとの残高・借方合計・貸方合計の登録
\def\register@Acc@Mon@Bal{%
  \ifnum\BKtemp@countA<\BK@closingmonth
    \global\advance\BKtemp@countA\@ne\relax
    \csxdef{Account\reserved@a @Month\the\BKtemp@countA @Balance}{0}%
    \csxdef{Account\reserved@a @Month\the\BKtemp@countA @Debit}{0}%
    \csxdef{Account\reserved@a @Month\the\BKtemp@countA @Credit}{0}%
    \expandafter\register@Acc@Mon@Bal
  \fi
}
%    \end{macrocode}
% \end{macro}
% \begin{macro}{\Account@getAccID}\marg{勘定科目名}\par
%  勘定科目名から勘定科目ID\meta{AccID}を得る．
%    \begin{macrocode}
\def\Account@getAccID#1{\csuse{revAccount#1@ID}}
%    \end{macrocode}
% \end{macro}
% \begin{macro}{\Account@getAccType}\marg{勘定科目名}\par
%  勘定科目名から勘定科目の種類\meta{AccType}を得る．
%    \begin{macrocode}
\def\Account@getAccType#1{\csuse{revAccount#1@Type}}
%    \end{macrocode}
% \end{macro}
% \begin{macro}{\Account@getAccTypeAccID}\marg{勘定科目名}\par
%  勘定科目名から\meta{AccType}\meta{AccID}を得る．
%    \begin{macrocode}
\def\Account@getAccTypeAccID#1{\Account@getAccType{#1}\Account@getAccID{#1}}

%%%%%%%%%%%%%%%%%%%%%%%%%%%%%%%%%%%%%%%%%%%%%%%%%%%%%%%%%%%%%%%%%%%%%%%%%%%%%%%%
%    \end{macrocode}
% \end{macro}
% \subsection{仕訳と仕訳帳}
% \subsubsection{仕訳}
% 取引の仕訳．
% 取引仕訳の総登録数は\cntname{JECount}カウンタで得られる．
%
% 取引IDを\meta{JEID}で表す．
% 取引仕訳の各要素には\cntname{\cmd{\JE}\meta{JEID}@\meta{prop.}}で
% アクセスできる．
% \meta{prop.}の種類は次の通り．
% \begin{description}
% \item[\cntname{Month}]
%   取引月を得る．
% \item[\cntname{Day}]
%   取引日を得る．
% \item[\cntname{Debit}]
%   借方の各要素を得る．
%   \begin{description}
%   \item[\cntname{Account}]
%     勘定科目名．
%   \item[\cntname{Value}]
%     取引額．
%   \item[\cntname{LEID}]
%     総勘定元帳の登録ID（LEID）．
%   \end{description}
% \item[\cntname{Credit}]
%   貸方の各要素を得る．
%   \begin{description}
%   \item[\cntname{Account}]
%     勘定科目名．
%   \item[\cntname{Value}]
%     取引額．
%   \item[\cntname{LEID}]
%     総勘定元帳の登録ID（LEID）．
%   \end{description}
% \item[\cntname{AccountDesc}]
%   摘要を得る．
% \item[\cntname{Line}]
%   仕訳帳での行数．
% \end{description}
%
% また摘要の有無を判定するために
% \begin{description}
% \item[\cmd{\ifJE\meta{JEID}@AccountDesc}]
% \item[\cmd{\JE@IfExistAccountDesc}]\marg{JEID}\marg{true}\marg{false}
% \end{description}
% が用意されている．
%
% 現時点では単一仕訳方式のみに対応しているが，
% 頑張れば1:N複合仕訳方式には対応できそうである（N:Nは対応しないかも）．
% ただし頑張るかは未定．
\iffalse
取引の仕訳．
取引仕訳の総登録数はJECountカウンタで得られる．

取引IDを<JEID>と表すことにすると
* JE<JEID>@Month: 取引月
* JE<JEID>@Day: 取引日
* JE<JEID>@Debit: 借方
  * Account: 勘定科目名
  * Value: 取引額
  * LEID: LEID
* JE<JEID>@Credit: 貸方
  * Account
  * Value
  * LEID
* JE<JEID>@AccountDesc: 摘要
* ifJE<JEID>@AccountDesc: 摘要の有無
* \JE@IfExistAccountDesc{<JEID>}{<true>}{<false>}: 摘要の有無 (true or false)で分岐するIf関数
* JE<JEID>@Line: 仕訳帳での行数

複合仕訳にするとなると科目の個数と各科目についてID割り振りが必要かなぁ……
面倒なのでやらないと思う．
\fi
% \begin{macro}{\registerJE@Date}\par
%    \begin{macrocode}
%% journalentry (JE)
\def\registerJE@Date#1#2{%
  \csgdef{JE\the\c@JECount @Month}{#1}%
  \csgdef{JE\the\c@JECount @Day}{#2}%
}
%    \end{macrocode}
% \end{macro}
% \begin{macro}{\registerJE@Debit}\par
%    \begin{macrocode}
\def\registerJE@Debit#1#2{%
  \csgdef{JE\the\c@JECount @Debit@Accouunt}{#1}%
  \csgdef{JE\the\c@JECount @Debit@Value}{#2}%
}
%    \end{macrocode}
% \end{macro}
% \begin{macro}{\registerJE@Credit}\par
%    \begin{macrocode}
\def\registerJE@Credit#1#2{%
  \csgdef{JE\the\c@JECount @Credit@Accouunt}{#1}%
  \csgdef{JE\the\c@JECount @Credit@Value}{#2}%
}
%    \end{macrocode}
% \end{macro}
% \begin{macro}{\registerJE@AccountDesc}\par
%    \begin{macrocode}
\def\registerJE@AccountDesc#1#2{%
  \csgdef{JE\the\c@JECount @AccountDesc}{#1}%
  \csgdef{ifJE\the\c@JECount @AccountDesc}{#2}%
}
%    \end{macrocode}
% \end{macro}
% \begin{macro}{\journ@lentry}\par
%    \begin{macrocode}
\newcommand{\journ@lentry}[6]{%
  \stepcounter{JECount}%
  \registerJE@Date{#1}{#2}%
  \registerJE@Debit{#3}{#4}%
  \registerJE@Credit{#5}{#4}%
  \registerJE@AccountDesc{#6}{\iftrue}%
  \csgdef{JE\the\c@JECount @Line}{3}%
  \ledger@entry{#3}{#5}{D}{#4}%
  \ledger@entry{#5}{#3}{C}{#4}%
}
%    \end{macrocode}
% \end{macro}
% \begin{macro}{\sjourn@lentry}\par
%    \begin{macrocode}
\newcommand{\sjourn@lentry}[5]{%
  \stepcounter{JECount}%
  \registerJE@Date{#1}{#2}%
  \registerJE@Debit{#3}{#4}%
  \registerJE@Credit{#5}{#4}%
  \registerJE@AccountDesc{}{\iffalse}%
  \csgdef{JE\the\c@JECount @Line}{2}%
  \ledger@entry{#3}{#5}{D}{#4}%
  \ledger@entry{#5}{#3}{C}{#4}%
}
%    \end{macrocode}
% \end{macro}
% \begin{macro}{\JE@IfExistAccountDesc}\par
%    \begin{macrocode}
\def\JE@IfExistAccountDesc#1#2#3{%
  \csname ifJE#1@AccountDesc\endcsname
  #2\else#3\fi
}

%    \end{macrocode}
% \end{macro}
% \begin{macro}{\journalentry}\marg{月}%
%                             \marg{日}%
%                             \marg{借方科目名}%
%                             \marg{取引額}%
%                             \marg{貸方科目名}%
%                             \marg{摘要}%\par
% \begin{macro}{\journalentry*}\marg{月}%
%                             \marg{日}%
%                             \marg{借方科目名}%
%                             \marg{取引額}%
%                             \marg{貸方科目名}\par
% 星付きのときは摘要が不要なものを示す．
%    \begin{macrocode}
%% \journalentry
%% #1: month
%% #2: date
%% #3: debit account name
%% #4: debit/credit value
%% #5: credit account name
%%   and if not stared
%% #6: description
\newcommand{\journalentry}{%
  \@ifstar{\sjourn@lentry}{\journ@lentry}%
}
%%%%%%%%%%%%%%%%%%%%%%%%%%%%%%%%%%%%%%%%%%%%%%%%%%%%%%%%%%%%%%%%%%%%%%%%%%%%%%%%
%    \end{macrocode}
% \end{macro}
% \end{macro}
% \subsubsection{仕訳帳}
% 仕訳帳の出力．
% 各\meta{JEID}について書き出していくだけだが，
% 次ページ・前ページ繰り越しのために\cmd{\sum@debit}および\cmd{\sum@credit}の加算も行っている．
% \begin{macro}{\printJE@Date}\par
%    \begin{macrocode}
\def\printJE@Date#1{%
  \csuse{JE#1@Month}/\csuse{JE#1@Day}%
}
%    \end{macrocode}
% \end{macro}
% \begin{macro}{\PJ@x@printline}\par
%    \begin{macrocode}
\def\PJ@x@printline#1#2#3#4#5#6#7{%
  #1&#2&#3&#4&#5&#6&#7\\%
}
%    \end{macrocode}
% \end{macro}
% \begin{macro}{\PJ@s@printline}\par
%    \begin{macrocode}
\def\PJ@s@printline#1#2#3#4#5#6{%
  #1&#2&\multicolumn{2}{l|}{#3}&#4&#5&#6\\%
}
%    \end{macrocode}
% \end{macro}
% \begin{macro}{\PJ@printline}\par
%    \begin{macrocode}
\def\PJ@printline{%
  \@ifstar{\PJ@s@printline}{\PJ@x@printline}%
}
%    \end{macrocode}
% \end{macro}
% \begin{macro}{\PJ@Transaction}\par
%    \begin{macrocode}
\def\PJ@Transaction#1{% #1: journal entry ID
  \JE@IfExistAccountDesc{#1}{%
    \gdef\@tmp@PJ{%
      \PJ@printline*{}{}{\csuse{JE#1@AccountDesc}}{}{}{}\hline
    }%
  }{%
    \global\let\@tmp@PJ\hline
  }%
  \edef\@tmp@PJ@deb{\csuse{JE#1@Debit@LEID}}%
  \PJ@printline{#1\label{JEID:#1}}{\printJE@Date{#1}}
               {（\csname JE#1@Debit@Accouunt\endcsname）}{}
               {\hyperref[LEID:\csuse{JE#1@Debit@LEID}]%
                         {\csuse{JE#1@Debit@LEID}}}
               {\csname JE#1@Debit@Value\endcsname}{}%
  \edef\@tmp@PJ@cre{\csuse{JE#1@Credit@LEID}}%
  \PJ@printline{}{}
               {}{（\csname JE#1@Credit@Accouunt\endcsname）}
               {\hyperref[LEID:\csuse{JE#1@Credit@LEID}]%
                         {\csuse{JE#1@Credit@LEID}}}
               {}{\csname JE#1@Credit@Value\endcsname}%
  \@tmp@PJ
}
%    \end{macrocode}
% \end{macro}
% \begin{macro}{\PJ@CarryPrev}\par
%    \begin{macrocode}
\def\PJ@CarryPrev{\gdef\n@xtprev{}%
  \PJ@printline{}{}{前頁繰越}{}{}
               {\the\sum@debit}{\the\sum@credit}\hline
}
%    \end{macrocode}
% \end{macro}
% \begin{macro}{\PJ@CarryForward}\par
%    \begin{macrocode}
\def\PJ@CarryForward{%
  \PJ@printline{}{}{}{次頁繰越}{}
               {\the\sum@debit}{\the\sum@credit}\hline
}
%    \end{macrocode}
% \end{macro}
% \begin{macro}{\PJ@Terminate}\par
%    \begin{macrocode}
\def\PJ@Terminate{%
  \PJ@printline{}{}{}{総計}{}
               {\the\sum@debit}{\the\sum@credit}\hline
}
%    \end{macrocode}
% \end{macro}
% \begin{macro}{\PJ@Loop@Table}\par
%    \begin{macrocode}
\def\PJ@Loop@Table{%
  \gdef\n@xt{}%
  \ifnum\c@JEIDiterator>\numexpr\c@JECount-1\relax
    \PJ@Terminate
  \else
    \ifnum\c@tableLineCount>\numexpr\maxtablelinenumber-1\relax
      \gdef\n@xt{\PJ@Loop}%
      \gdef\n@xtprev{\PJ@CarryPrev}%
      \PJ@CarryForward
    \else
      \refstepcounter{JEIDiterator}%
      \global\advance\c@tableLineCount
        \csuse{JE\the\c@JEIDiterator @Line}\relax
      \global\advance\sum@debit
        \csuse{JE\the\c@JEIDiterator @Debit@Value}\relax
      \global\advance\sum@credit
        \csuse{JE\the\c@JEIDiterator @Credit@Value}\relax
      \PJ@Transaction{\the\c@JEIDiterator}%
      \expandafter\expandafter\expandafter\PJ@Loop@Table\expandafter\expandafter
%     \gdef\n@xt{\PJ@Loop@Table}%
    \fi
  \fi
}
%    \end{macrocode}
% \end{macro}
% \begin{macro}{\PJ@Loop}\par
%    \begin{macrocode}
\def\PJ@Loop{%
  \begin{table}[ht]\centering
  \begin{tabular}{r|D{/}{/}{2/2}|lr|r|r|r}
    \makebox[2\zw]{ID}%
    & \multicolumn{1}{c|}{date}%
    & \makebox[10\zw][c]{摘}
    & \makebox[10\zw][c]{要}%
    & \makebox[2\zw][c]{LEID}%
    & \makebox[5\zw][c]{借方}%
    & \makebox[5\zw][c]{貸方} \\ \hline
    \n@xtprev
    \PJ@Loop@Table
  \end{tabular}
  \end{table}
  \clearpage
  \global\c@tableLineCount\z@\relax
  \n@xt
}
%    \end{macrocode}
% \end{macro}
% \begin{macro}{\printjournal}\par
%    \begin{macrocode}
\newcommand{\printjournal}{%
  \setcounter{JEIDiterator}{0}%
  \gdef\n@xtprev{}%
  \gdef\n@xt{}%
  \PJ@Loop
}
%%%%%%%%%%%%%%%%%%%%%%%%%%%%%%%%%%%%%%%%%%%%%%%%%%%%%%%%%%%%%%%%%%%%%%%%%%%%%%%%
%    \end{macrocode}
% \end{macro}
% \subsection{総勘定元帳}
% \subsubsection{総勘定元帳への転記登録}
% 仕訳の登録と同時に総勘定元帳への登録も行う．
% 
% 勘定科目のIDを\meta{AccType}\meta{AccID}，
% 各勘定科目ごとの取引登録IDを\meta{AccLdgID}で表すこととする．
% 各勘定科目ごとの取引の各要素には
% \cntname{\cmd{\Account}\meta{AccType}\meta{AccID}@Ledger\meta{AccLdgID}@\meta{prop.}}で
% アクセスできる．
% \meta{prop.}の種類は次の通り．
% \begin{description}
% \item[\cntname{FullLEID}]
%   登録された取引の描画用の形式に直されたLEID
%   （\cntname{\meta{AccType}\meta{AccID}-\meta{AccLdgID}}）を得る．
% \item[\cntname{JEID}]
%   登録された取引の取引IDを得る．
% \item[\cntname{PartnerAccount}]
%   登録された取引の相手勘定科目を得る．
% \item[\cntname{JEDesc}]
%   登録された取引の摘要を得る．
% \item[\cntname{IsAccountDebit}]
%   登録された取引での借貸がどちらだったか（\cntname{D}または\cntname{C}）を得る．
% \item[\cntname{AccountValue}]
%   登録された取引での取引額を得る．
% \end{description}
% 
\iffalse
総勘定元帳．
勘定科目ID(AccID)および各勘定科目ごとの取引登録ID(AccLdgID)，
それによって参照される要素を持つ．

!externな内容
* Account<AccType><AccID>@LedgerCount: 勘定科目に登録された取引登録数．
                                        この数字でイテレータの分岐をさせる．
!externな内容ここまで

* Account<AccType><AccID>@Ledger<AccLdgID>@FullLEID: 描画用のLEID<AccType><AccID>-<AccLdgID>を得る．
* Account<AccType><AccID>@Ledger<AccLdgID>@JEID: 登録された取引の取引ID．
* Account<AccType><AccID>@Ledger<AccLdgID>@PartnerAccount: 登録された取引の相手勘定科目．
* Account<AccType><AccID>@Ledger<AccLdgID>@JEDesc: 登録された取引の摘要．
* Account<AccType><AccID>@Ledger<AccLdgID>@IsAccountDebit: 登録された取引での借貸がどちらだったか．
* Account<AccType><AccID>@Ledger<AccLdgID>@AccountValue: 登録された取引での取引額．

JEIDで取引の月日が分かるのでデータとして持つ必要はない．
一方で相手科目は持たせた方が楽っぽい？
必ずしも持たせる必要はないと思うけど処理がめんどいと思う．
\fi
% \begin{macro}{\ledger@entry}\marg{勘定科目名}%
%                             \marg{相手科目名}%
%                             \marg{借（貸）方}%
%                             \marg{取引額}\par
% \cntname{\cmd{\Account}\meta{AccType}\meta{AccID}@Ledger\meta{AccLdgID}@\meta{prop.}}に
% 各種値を代入していく．
% そのほか\cntname{\cmd{\Account}\meta{AccType}\meta{AccID}@Month\meta{Month}@\meta{prop.}}に
% 各種値を代入していく．
% \meta{借（貸）方}はその取引で借方のときに\cntname{D}を，貸方の時に\cntname{C}をとる．
% また\meta{取引額}は\cntname{0}のときに月次締切を表すことにしている\footnote{%
%   \meta{借（貸）方}にフラグを持たせるほうがいいのかもしれない．
% }．
%    \begin{macrocode}
%% general ledger (GL)
%% #1: Account name
%% #2: Partner account name
%% #3: is account Debit? (boolean; D or C)
%% #4: Account value(if this is 0 then this ledger entry is Monthly closing)
\def\ledger@entry#1#2#3#4{%
  \edef\LE@tmpA{\Account@getAccTypeAccID{#1}}% \LE@tmpA <- <AccType><AccID>
  \BKtemp@countB=\csuse{Account\LE@tmpA @LedgerCount}\relax
  % tempcountB <- Account<AccType><AccID>@LedgerCount
  \advance\BKtemp@countB\@ne
  \csxdef{Account\LE@tmpA @LedgerCount}{\the\BKtemp@countB}%
  \csxdef{Account\LE@tmpA @Ledger\the\BKtemp@countB @FullLEID}%
    {\LE@tmpA-\the\BKtemp@countB}%
  \csxdef{Account\LE@tmpA @Ledger\the\BKtemp@countB @JEID}{%
    \the\c@JECount}% refference JEID
  \csxdef{Account\LE@tmpA @Ledger\the\BKtemp@countB @PartnerAccount}{#2}%
  \csxdef{Account\LE@tmpA @Ledger\the\BKtemp@countB @JEDesc}%
    {\csuse{JE\the\c@JECount @AccountDesc}}%
  \csgdef{Account\LE@tmpA @Ledger\the\BKtemp@countB @IsAccountDebit}{#3}%
  \csgdef{Account\LE@tmpA @Ledger\the\BKtemp@countB @AccountValue}{#4}%
  % Account balance, Debit and Credit in the month
  \BKtemp@countA=\csuse{JE\the\c@JECount @Month}\relax% tempcountA <- JE@Month
  \csxdef{Account\LE@tmpA @Month\the\BKtemp@countA @Balance}{%
    \the\numexpr\csuse{Account\LE@tmpA @Month\the\BKtemp@countA @Balance}%
    + \csuse{Account\LE@tmpA @Ledger\the\BKtemp@countB @AccountValue}\relax
  }%
  \if#3D%
    \csxdef{Account\LE@tmpA @Month\the\BKtemp@countA @Debit}{%
      \the\numexpr\csuse{Account\LE@tmpA @Month\the\BKtemp@countA @Debit}%
      + \csuse{Account\LE@tmpA @Ledger\the\BKtemp@countB @AccountValue}\relax
    }%
    \csxdef{JE\csuse{Account\LE@tmpA @Ledger\the\BKtemp@countB @JEID}@Debit@LEID}{%
      \csuse{Account\LE@tmpA @Ledger\the\BKtemp@countB @FullLEID}%
    }%
  \fi
  \if#3C%
    \csxdef{Account\LE@tmpA @Month\the\BKtemp@countA @Credit}{%
      \the\numexpr\csuse{Account\LE@tmpA @Month\the\BKtemp@countA @Credit}%
      + \csuse{Account\LE@tmpA @Ledger\the\BKtemp@countB @AccountValue}\relax
    }%
    \csxdef{JE\csuse{Account\LE@tmpA @Ledger\the\BKtemp@countB @JEID}@Credit@LEID}{%
      \csuse{Account\LE@tmpA @Ledger\the\BKtemp@countB @FullLEID}%
    }%
  \fi
}
%    \end{macrocode}
% \end{macro}
% \begin{macro}{\ledgermonthlyclosing}\marg{締切月}\par
% 月次締切用コマンド．
% 各\meta{AccType}および各\meta{AccID}についてイテレータ
% \cmd{\BKtemp@countD}および\cmd{\BKtemp@countC}を用いてループを回す．
%    \begin{macrocode}
\newcommand{\ledgermonthlyclosing}[1]{% #1: closing month
  \BKtemp@countC\z@\relax
  \BKtemp@countD\z@\relax
  \stepcounter{ClosingMonthCount}%
  \gdef\LMC@month{#1}%
  \ledgermonthlyclosing@loop
}
%    \end{macrocode}
% \end{macro}
% \begin{macro}{\ledgermonthlyclosing@loop}\par
% \changes{v1.1.0}{2024/02/27}
%   {勘定科目の種類についてリファクタしたのに伴ってアクセス用のマクロを変更．}
%    \begin{macrocode}
\def\ledgermonthlyclosing@loop{% loop for each AccType
  \ifnum\the\BKtemp@countD<\c@AccountTypeCount\relax
    \global\advance\BKtemp@countD\@ne
    \global\BKtemp@countC\z@\relax
    \xdef\LMC@AccType{\AccType@getAccType{\the\BKtemp@countD}}%
    % \LMC@AccType <- AccType
    \expandafter\ledgermonthlyclosing@l@@p
  \fi
}
%    \end{macrocode}
% \end{macro}
% \begin{macro}{\ledgermonthlyclosing@l@@p}\par
% \changes{v1.1.0}{2024/02/27}
%   {勘定科目の種類についてリファクタしたのに伴ってアクセス用のCSを変更．}
%    \begin{macrocode}
\def\ledgermonthlyclosing@l@@p{% loop for each AccID
  \ifnum\BKtemp@countC<\csname c@AccType\LMC@AccType @AccountCount\endcsname
    \global\advance\BKtemp@countC\@ne
    \xdef\reserved@a{\csuse{Account\LMC@AccType\the\BKtemp@countC @Name}}%
    \expandafter\ledger@entry\expandafter{\reserved@a}{\LMC@month}{\relax}{0}%
    \expandafter\ledgermonthlyclosing@l@@p
  \else
    \expandafter\ledgermonthlyclosing@loop
  \fi
}
%    \end{macrocode}
% \end{macro}
% \begin{macro}{\ledgerclosing}\marg{損益勘定科目名}\marg{利益剰余金勘定科目名}\par
% 年度締切用コマンド（まだ未実装）．
% \begin{enumerate}
% \item 収益と費用の各勘定科目について損益振替をする
% \item 総資産の勘定科目（一般には利益剰余金勘定）に損益の振替をする
% \item 資産，負債および総資産の勘定科目について閉鎖残高に振替をする
% \end{enumerate}
% が目標。
%    \begin{macrocode}
\newcommand{\ledgerclosing}{%
}
%    \end{macrocode}
% \end{macro}
% \begin{macro}{\journalentry@PLtransfar}\marg{損益勘定科目名}\par
% \changes{v1.1.0}{2024/02/19}
%   {ADD：損益振替仕訳の準備}%
% 収益と費用の各勘定科目について損益振替をする．
% \cmd{\PLtransfar@AccTypeAccID}に損益勘定科目の\meta{AccType}\meta{AccID}を
% 控えておく．
%    \begin{macrocode}
\newcommand{\journalentry@PLtransfar}[1]{%
  \edef\PLtransfar@AccTypeAccID{\Account@getAccTypeAccID{#1}}%
  \setcounter{AccountIDiterator}{0}%
  \journalentry@PLtransfar@LoopAccType{R}% 収益 Revenue
  \journalentry@PLtransfar@LoopAccType{E}% 費用 Expence
}
%    \end{macrocode}
% \end{macro}
% \begin{macro}{\journalentry@PLtransfar@LoopAccType}\marg{勘定科目の種類}\par
% \changes{v1.1.0}{2024/02/19}
%   {ADD：損益振替仕訳の準備}%
% 指定した勘定科目の種類に属する勘定科目について処理をする．
%    \begin{macrocode}
\newcommand{\journalentry@PLtransfar@LoopAccType}[1]{%
  \ifnum\c@AccountIDiterator<\csuse{c@AccType#1@AccountCount}%
    \stepcounter{AccountIDiterator}%
    \xdef\JE@PLT@Account{\csuse{Account#1\the\c@AccountIDiterator@Name}}%
    \journalentry@PLtransfar@register
    \expandafter\@firstonone
  \else
    \expandafter\@gobble
  \fi{\journalentry@PLtransfar@LoopAccType{#1}}%
}
%    \end{macrocode}
% \end{macro}
% \begin{macro}{\journalentry@PLtransfar@register}\par
% \changes{v1.1.0}{2024/02/19}
%   {ADD：損益振替仕訳の準備}%
% 指定した勘定科目の種類に属する勘定科目について処理をする．
% ……んですが，残高をどこから得ればいいかの設計で迷い中．
% 貸借の分岐は勘定科目の要素\meta{LedgerBakanceType}で行うのがスマートだろうか．
%    \begin{macrocode}
\newcommand{\journalentry@PLtransfar@register}{%
%%\journalentry*{13}{01}{\JE@PLT@Account}
}
%    \end{macrocode}
% \end{macro}
% \subsubsection{総勘定元帳の出力}
% 一貫して\cmd{\PGL@}を接頭辞として用いるようにする．
% \begin{macro}{\PGL@printline}%
%   \marg{AccLdgID}%
%   \marg{月/日}%
%   \marg{相手勘定科目名}%
%   \marg{摘要}%
%   \marg{JEID}%
%   \marg{借方}%
%   \marg{貸方}%
%   \marg{残高賃借}%
%   \marg{残高}\par
% 総勘定元帳での各取引の出力用．
%    \begin{macrocode}
% #1: LEID
% #2: Month/Day
% #3: Account title
% #4: Desc.
% #5: JEID
% #6: Debit
% #7: Credit
% #8: Balance type
% #9: Balance
\def\PGL@printline#1#2#3#4#5#6#7#8#9{%
  #1&#2&#3&#4&#5&#6&#7&#8&#9\\%
}

%    \end{macrocode}
% \end{macro}
% \begin{macro}{\PGL@Transaction}\marg{AccLdgID}\par
% \changes{v1.1.0}{2024/02/27}
%   {Account Typeを一時的に保持するマクロの名前をよりユニークな名前に変更．}
% \meta{AccLdgID}の取引を出力する．
%    \begin{macrocode}
\def\PGL@Transaction#1{% #1: ladger entry ID
  \xdef\PGL@tmp@AccTypeAccID{\PGL@tmp@AccType\the\c@AccountIDiterator}%
  \global\BKtemp@countB=%
    \csuse{Account\PGL@tmp@AccTypeAccID @Ledger#1@JEID}\relax%
    % tempcountB <- Account[AccType][AccID].Ledger[AccLdgID].JEID
  \xdef\Acc@Ledg@IsAccDeb{%
    \csuse{Account\PGL@tmp@AccTypeAccID @Ledger#1@IsAccountDebit}%
  }%
  \xdef\Acc@LedgBal{%
    \csuse{Account\PGL@tmp@AccTypeAccID @LedgerBalanceType}%
  }%
  \PGL@printline{% LEID
    \csuse{Account\PGL@tmp@AccTypeAccID @Ledger#1@FullLEID}%
    \label{LEID:\PGL@tmp@AccTypeAccID-#1}%
  }{% Month/Day
    \printJE@Date{\the\BKtemp@countB}%
  }{% Account title
    \csuse{Account\PGL@tmp@AccTypeAccID @Ledger#1@PartnerAccount}%
  }{% Desc.
    \csuse{Account\PGL@tmp@AccTypeAccID @Ledger#1@JEDesc}%
  }{% JEID
    \hyperref[JEID:\the\BKtemp@countB]{\the\BKtemp@countB}%
  }{% Debit
    \if\Acc@Ledg@IsAccDeb D\relax%
      \csuse{Account\PGL@tmp@AccTypeAccID @Ledger#1@AccountValue}%
      \global\advance\sum@debit
                     \csuse{Account\PGL@tmp@AccTypeAccID
                             @Ledger#1@AccountValue}\relax
    \fi
  }{% Credit
    \if\Acc@Ledg@IsAccDeb C\relax%
      \csuse{Account\PGL@tmp@AccTypeAccID @Ledger#1@AccountValue}%
      \global\advance\sum@credit
                     \csuse{Account\PGL@tmp@AccTypeAccID
                             @Ledger#1@AccountValue}\relax
    \fi
  }{% Balance type
    \if\Acc@LedgBal D%
      借\else 貸%
    \fi
  }{% Balance
    \if\Acc@Ledg@IsAccDeb\Acc@LedgBal
      \global\advance\c@LedgerEntryBalance
                     \csuse{Account\PGL@tmp@AccTypeAccID
                             @Ledger#1@AccountValue}\relax
    \else
      \global\advance\c@LedgerEntryBalance
                    -\csuse{Account\PGL@tmp@AccTypeAccID
                             @Ledger#1@AccountValue}\relax
    \fi \the\c@LedgerEntryBalance}%
}
%    \end{macrocode}
% \end{macro}
\iffalse
  \printgeneralledger
  v
  \PGL@Loop -> exit loop
  v
  v
  \PGL@L@@p -> \PGL@Loop
  v
  \PGL@ACLoop
  v
  \PGL@ACLoop@Table -> \PGL@CarryForward -> \PGL@ACLoop
  v ^
  v ^               \> \PGL@Terminate    -> \PGL@L@@p
  v ^
  \PGL@Transaction
\fi
% \begin{macro}{\printgeneralledger}\par
% 実際に総勘定元帳を出力するコマンド．
% いくつかのパラメタを初期化して\cmd{\PLG@Loop}を呼び出す．
%    \begin{macrocode}
\newcommand{\printgeneralledger}{%
  \setcounter{AccountTypeiterator}{0}%
  \gdef\n@xtprev{}%
  \gdef\n@xt{}%
  \PGL@Loop
}
%    \end{macrocode}
% \end{macro}
% \begin{macro}{\PGL@Loop}\par
% \changes{v1.1.0}{2024/02/27}
%   {勘定科目の種類についてリファクタしたのに伴ってアクセス用のマクロを変更．}
% \changes{v1.1.0}{2024/02/27}
%   {Account Typeを一時的に保持するマクロの名前をよりユニークな名前に変更．}
% 各勘定科目の種類について\cmd{PLG@L@@p}を呼び出す．
%    \begin{macrocode}
\def\PGL@Loop{%
  \ifnum\c@AccountTypeiterator<\c@AccountTypeCount\relax
    \stepcounter{AccountTypeiterator}%
    \xdef\PGL@tmp@AccType{\AccType@getAccType{\the\c@AccountTypeiterator}}%
    \setcounter{AccountIDiterator}{0}%
    \expandafter\PGL@L@@p
  \fi
}
%    \end{macrocode}
% \end{macro}
% \begin{macro}{\PGL@L@@p}\par
% \changes{v1.1.0}{2024/02/27}
%   {Account Typeを一時的に保持するマクロの名前をよりユニークな名前に変更．}
% \changes{v1.1.0}{2024/02/27}
%   {勘定科目の種類についてリファクタしたのに伴ってアクセス用のCSを変更．}
% ある勘定科目の種類に属する各勘定科目について\cmd{PLG@ACLoop}を呼び出す．
% 全ての勘定科目について呼び出したら\cmd{PLG@Loop}に処理を戻す．
%    \begin{macrocode}
\def\PGL@L@@p{%
  \ifnum\c@AccountIDiterator<\csname c@AccType\PGL@tmp@AccType @AccountCount\endcsname
    \stepcounter{AccountIDiterator}%
    \expandafter\subsection\expandafter%
      {\csuse{Account\PGL@tmp@AccType\the\c@AccountIDiterator @Name}}%
    \setcounter{AccLdgIDiterator}{0}%
    \setcounter{LedgerEntryBalance}{0}%
    \global\sum@debit\z@\relax\global\sum@credit\z@\relax
    \expandafter\PGL@ACLoop
  \else
    \expandafter\PGL@Loop
  \fi
}
%    \end{macrocode}
% \end{macro}
% \begin{macro}{\PGL@ACLoop}\par
% 勘定元帳の表を作り，（改ページが
% 1ページ当たりの行数を超えたことに因るときには前頁繰越をして）
% \cmd{PLG@ACLoop@Table}を呼び出す．
% その処理が終わった後改ページをして\cmd{PLG@ACLoop@Table}内で設定した
% \cmd{n@xt}を呼び出す：これは改ページが
% 1ページ当たりの行数を超えたことに因るときには\cmd{PGL@ACLoop}，
% 勘定科目の仕訳を全て出力したことに因るときには\cmd{PGL@L@@p}である．
% 
%    \begin{macrocode}
\def\PGL@ACLoop{%
  \begin{table}[ht]\centering
  \begin{tabular}{r|D{/}{/}{2/2}|l|l|r|r|r|c|r}
  \PGL@printline{\makebox[4\zw]{LEID}}%
                {\multicolumn{1}{c|}{\makebox[3\zw]{date}}}%
                {\makebox[7\zw]{相手勘定科目}}%
                {\multicolumn{1}{c|}{\makebox[10\zw]{摘}\makebox[10\zw]{要}}}%
                {\makebox[2\zw]{JEID}}%
                {\makebox[4\zw]{借方}}%
                {\makebox[4\zw]{貸方}}%
                {\makebox[2\zw]{借/貸}}%
                {\makebox[4\zw]{残高}}\hline
    \n@xtprev
    \PGL@ACLoop@Table
  \end{tabular}
  \end{table}
  \clearpage
  \global\c@tableLineCount\z@\relax
  \n@xt
}
%    \end{macrocode}
% \end{macro}
% \begin{macro}{\PGL@ACLoop@Table}\par
% \changes{v1.1.0}{2024/02/27}
%   {Account Typeを一時的に保持するマクロの名前をよりユニークな名前に変更．}
% ある勘定科目の各仕訳について\cmd{PGL@Transaction}を呼び出す．
% 全ての仕訳を出力したら\cmd{PGL@Terminate}を呼び出す．
% その他改ページ処理など．
%    \begin{macrocode}
\def\PGL@ACLoop@Table{%
  \ifnum\c@AccLdgIDiterator>\numexpr
      \csuse{Account\PGL@tmp@AccType\the\c@AccountIDiterator @LedgerCount}-1\relax
    \hline
    \PGL@Terminate
  \else
    \ifnum\c@tableLineCount>\numexpr\maxtablelinenumber-1\relax
      \hline
      \gdef\n@xt{\PGL@ACLoop}%
      \gdef\n@xtprev{\PGL@CarryPrev}%
      \PGL@CarryForward
    \else
      \gdef\n@xt{\PGL@L@@p}%
      \stepcounter{AccLdgIDiterator}%
      \ifnum\csname Account\PGL@tmp@AccType\the\c@AccountIDiterator
                    @Ledger\the\c@AccLdgIDiterator
                    @AccountValue\endcsname=0\relax
        \stepcounter{tableLineCount}%
        \PGL@MonthlyClosing%
      \else
        \stepcounter{tableLineCount}%
        \PGL@Transaction{\the\c@AccLdgIDiterator}%
      \fi
      \expandafter\expandafter\expandafter\PGL@ACLoop@Table
      \expandafter\expandafter
    \fi
  \fi
}
%    \end{macrocode}
% \end{macro}
% \begin{macro}{\PGL@CarryPrev}\par
%    \begin{macrocode}
\def\PGL@CarryPrev{%
  \gdef\n@xtprev{}%
  \PGL@printline
    {}{}{前頁繰越}{}{}%
    {\the\sum@debit}{\the\sum@credit}{}{\the\c@LedgerEntryBalance}%
  \hline
}
%    \end{macrocode}
% \end{macro}
% \begin{macro}{\PGL@MonthlyClosing}\par
% \changes{v1.1.0}{2024/02/27}
%   {Account Typeを一時的に保持するマクロの名前をよりユニークな名前に変更．}
%    \begin{macrocode}
\def\PGL@MonthlyClosing{%
  \xdef\PGL@tmp@AccTypeAccID{\PGL@tmp@AccType\the\c@AccountIDiterator}%
  \xdef\Acc@LedgBal{%
    \csuse{Account\PGL@tmp@AccTypeAccID
           @LedgerBalanceType}}%
  \xdef\tmp@PGL@ClosingMonth{%
    \csuse{Account\PGL@tmp@AccTypeAccID
           @Ledger\the\c@AccLdgIDiterator
           @PartnerAccount}}% \tmp@PGL@ClosingMonth <- Closing month
  \global\sum@debit\z@\relax \global\sum@credit\z@\relax
  \csxdef{Account\PGL@tmp@AccTypeAccID
          @Month\tmp@PGL@ClosingMonth @Balance}{\the\c@LedgerEntryBalance}%
  \PGL@printline{}{}{\tmp@PGL@ClosingMonth 月決算}{}{}%
     {\csuse{Account\PGL@tmp@AccTypeAccID
             @Month\tmp@PGL@ClosingMonth @Debit}}%
     {\csuse{Account\PGL@tmp@AccTypeAccID
             @Month\tmp@PGL@ClosingMonth @Credit}}%
     {}{}\hline
}
%    \end{macrocode}
% \end{macro}
% \begin{macro}{\PGL@CarryForward}\par
%    \begin{macrocode}
\def\PGL@CarryForward{%
  \PGL@printline{}{}{\multicolumn{1}{r|}{次頁繰越}}{}{}%
                {\the\sum@debit}{\the\sum@credit}{}{\the\c@LedgerEntryBalance}%
  \hline
}
%    \end{macrocode}
% \end{macro}
% \begin{macro}{\PGL@Terminate}\par
%    \begin{macrocode}
\def\PGL@Terminate{%
  \PGL@printline{}{}{}{}{}{\the\sum@debit}{\the\sum@credit}{}{}\hline\hline
}

%%%%%%%%%%%%%%%%%%%%%%%%%%%%%%%%%%%%%%%%%%%%%%%%%%%%%%%%%%%%%%%%%%%%%%%%%%%%%%%%
%    \end{macrocode}
% \end{macro}
% \subsection{合計残高試算表}
% \begin{macro}{\TB@printTable}\par
%    \begin{macrocode}
% Trial balance (TB)
\def\TB@printTable#1#2#3#4#5{%
  #1&#2&#3&#4&#5\\%
}
%    \end{macrocode}
% \end{macro}
% \begin{macro}{\TB@L@@p@T@ble}\par
% \changes{v1.1.0}{2024/02/27}
%   {Account Typeを一時的に保持するマクロの名前をよりユニークな名前に変更．}
%    \begin{macrocode}
\def\TB@L@@p@T@ble{%
  \xdef\reserved@a{\TB@tmp@AccType\the\c@AccountIDiterator}%
  \xdef\Acc@LedgBal{%
    \csname Account\reserved@a @LedgerBalanceType\endcsname}%
  \TB@printTable{%
    \if\Acc@LedgBal D%
      \csuse{Account\reserved@a @Month\the\c@MonthIterator @Balance}%
      \global\advance\BKtemp@countA
                     \csuse{Account\reserved@a @Month\the\c@MonthIterator @Balance}\relax
      \csxdef{Account\TB@tmp@AccType @Month\the\c@MonthIterator @sumDebit}{%
        \the\numexpr\csuse{Account\TB@tmp@AccType @Month\the\c@MonthIterator @sumDebit}%
        + \csuse{Account\reserved@a @Month\the\c@MonthIterator @Balance}\relax
      }%
    \fi
  }{%
    \csuse{Account\reserved@a @Month\the\c@MonthIterator @Debit}%
    \global\advance\BKtemp@countB
                   \csuse{Account\reserved@a @Month\the\c@MonthIterator @Debit}\relax
    \csxdef{Account\TB@tmp@AccType @Month\the\c@MonthIterator @Debit}{%
      \the\numexpr\csuse{Account\TB@tmp@AccType @Month\the\c@MonthIterator @Debit}%
      + \csuse{Account\reserved@a @Month\the\c@MonthIterator @Debit}\relax
    }%
  }{%
    \csuse{Account\reserved@a @Name}%
  }{%
    \csuse{Account\reserved@a @Month\the\c@MonthIterator @Credit}%
    \global\advance\BKtemp@countC
                   \csuse{Account\reserved@a @Month\the\c@MonthIterator @Credit}\relax
    \csxdef{Account\TB@tmp@AccType @Month\the\c@MonthIterator @Credit}{%
      \the\numexpr\csuse{Account\TB@tmp@AccType @Month\the\c@MonthIterator @Credit}%
      + \csuse{Account\reserved@a @Month\the\c@MonthIterator @Credit}\relax
    }%
  }{%
    \if\Acc@LedgBal C%
      \csuse{Account\reserved@a @Month\the\c@MonthIterator @Balance}%
      \global\advance\BKtemp@countD
                     \csuse{Account\reserved@a @Month\the\c@MonthIterator @Balance}\relax
      \csxdef{Account\TB@tmp@AccType @Month\the\c@MonthIterator @sumCredit}{%
        \the\numexpr\csuse{Account\TB@tmp@AccType @Month\the\c@MonthIterator @sumCredit}%
        + \csuse{Account\reserved@a @Month\the\c@MonthIterator @Balance}\relax
      }%
    \fi
  }%
}
%    \end{macrocode}
% \end{macro}
% \begin{macro}{\printmonthlytrialbalance}\par
%    \begin{macrocode}
\newcommand{\printmonthlytrialbalance}{%
  \setcounter{MonthIterator}{0}%
  \gdef\n@xtprev{}%
  \gdef\n@xt{}%
% \typeout{1}%
  \TB@Loop
}
%    \end{macrocode}
% \end{macro}
% \begin{macro}{\TB@Loop}\par
%    \begin{macrocode}
\def\TB@Loop{%
  \ifnum\c@MonthIterator<\c@ClosingMonthCount
    \stepcounter{MonthIterator}%
    \BKtemp@countA\z@ \BKtemp@countB\z@
    \BKtemp@countC\z@ \BKtemp@countD\z@
    \expandafter\subsection\expandafter{\the\c@MonthIterator 月}%
    \setcounter{AccountTypeiterator}{0}%
%   \typeout{2: MonthIterator=\the\c@MonthIterator}%
    \expandafter\TB@L@@p
  \fi
}
%    \end{macrocode}
% \end{macro}
% \begin{macro}{\TB@L@@p}\par
%    \begin{macrocode}
\def\TB@L@@p{%
  \begin{table}[ht]\centering
  \begin{tabular}{rrcrr}
    \multicolumn{2}{c}{借方}&&\multicolumn{2}{c}{貸方}\\%
    \TB@printTable{\makebox[5\zw]{残高}}{\makebox[5\zw]{合計}}%
                  {\makebox[10\zw]{勘定科目}}{\makebox[5\zw]{合計}}%
                  {\makebox[5\zw]{残高}}\hline%
    \TB@Loop@Table
    \hline
    \TB@printTable@sumAccType
    \hline
    \TB@printTable{\the\BKtemp@countA}{\the\BKtemp@countB}%
                  {合計}%
                  {\the\BKtemp@countC}{\the\BKtemp@countD}\hline%
  \end{tabular}
  \end{table}
  \clearpage
  \TB@Loop
}
%    \end{macrocode}
% \end{macro}
% \begin{macro}{\TB@Loop@Table}\par
% \changes{v1.1.0}{2024/02/27}
%   {Account Typeを一時的に保持するマクロの名前をよりユニークな名前に変更．}
%    \begin{macrocode}
\def\TB@Loop@Table{%
  \ifnum\c@AccountTypeiterator<\c@AccountTypeCount\relax
    \stepcounter{AccountTypeiterator}%
    \xdef\TB@tmp@AccType{\AccType@getAccType{\the\c@AccountTypeiterator}}%
    \csxdef{Account\TB@tmp@AccType @Month\the\c@MonthIterator @sumDebit}{0}%
    \csxdef{Account\TB@tmp@AccType @Month\the\c@MonthIterator @Debit}{0}%
    \csxdef{Account\TB@tmp@AccType @Month\the\c@MonthIterator @sumCredit}{0}%
    \csxdef{Account\TB@tmp@AccType @Month\the\c@MonthIterator @Credit}{0}%
    \setcounter{AccountIDiterator}{0}%
%   \typeout{3: AccountTypeiterator=\TB@tmp@AccType}%
    \expandafter\TB@L@@p@Table
  \fi
}
%    \end{macrocode}
% \end{macro}
% \begin{macro}{\TB@L@@p@Table}\par
% \changes{v1.1.0}{2024/02/27}
%   {Account Typeを一時的に保持するマクロの名前をよりユニークな名前に変更．}
% \changes{v1.1.0}{2024/02/27}
%   {勘定科目の種類についてリファクタしたのに伴ってアクセス用のCSを変更．}
%    \begin{macrocode}
\def\TB@L@@p@Table{%
  \ifnum\c@AccountIDiterator<\csname c@AccType\TB@tmp@AccType @AccountCount\endcsname
    \stepcounter{AccountIDiterator}%
%   \typeout{4: AccountIDiterator=\the\c@AccountIDiterator}%
    \TB@L@@p@T@ble
    \expandafter\TB@L@@p@Table
  \else
    \TB@Loop@Table
  \fi
}
%    \end{macrocode}
% \end{macro}
% \begin{macro}{\TB@printTable@sumAccType}\par
%    \begin{macrocode}
\def\TB@printTable@sumAccType{%
  \TB@printTable{\csuse{AccountA@Month\the\c@MonthIterator @sumDebit}}%
                {\csuse{AccountA@Month\the\c@MonthIterator @Debit}}%
                {資産}%
                {\csuse{AccountA@Month\the\c@MonthIterator @Credit}}%
                {}%
  \TB@printTable{}%
                {\csuse{AccountL@Month\the\c@MonthIterator @Debit}}%
                {負債}%
                {\csuse{AccountL@Month\the\c@MonthIterator @Credit}}%
                {\csuse{AccountL@Month\the\c@MonthIterator @sumCredit}}%
  \TB@printTable{}%
                {\csuse{AccountN@Month\the\c@MonthIterator @Debit}}%
                {総資産}%
                {\csuse{AccountN@Month\the\c@MonthIterator @Credit}}%
                {\csuse{AccountN@Month\the\c@MonthIterator @sumCredit}}%
  \TB@printTable{}%
                {\csuse{AccountR@Month\the\c@MonthIterator @Debit}}%
                {収益}%
                {\csuse{AccountR@Month\the\c@MonthIterator @Credit}}%
                {\csuse{AccountR@Month\the\c@MonthIterator @sumCredit}}%
  \TB@printTable{\csuse{AccountE@Month\the\c@MonthIterator @sumDebit}}%
                {\csuse{AccountE@Month\the\c@MonthIterator @Debit}}%
                {費用}%
                {\csuse{AccountE@Month\the\c@MonthIterator @Credit}}%
                {}%
  \TB@printTable{\csuse{AccountM@Month\the\c@MonthIterator @sumDebit}}%
                {\csuse{AccountM@Month\the\c@MonthIterator @Debit}}%
                {その他}%
                {\csuse{AccountM@Month\the\c@MonthIterator @Credit}}%
                {\csuse{AccountM@Month\the\c@MonthIterator @sumCredit}}%
}
%    \end{macrocode}
% \end{macro}
%
% 以上で終わりです.
%
%    \begin{macrocode}
%</style>
\endinput
%    \end{macrocode}
%
%
% \Finale
% \PrintIndex \PrintChanges
